%======================================================================
\chapter{The Base Protocol}
\label{ch:proto}
%======================================================================

This chapter sets out the protocol exactly as proposed by Rupa \emph{et al.}
\cite{basepaper}: the network it assumes, its three phases, all twelve messages, and a
worked example following one real message from the seabed to the shore.

\section{Network model}

The network is a five-layer hierarchy. Data flows upward from the seabed to land, and every
adjacent pair authenticates each other before any payload moves.

\begin{itemize}
\item \textbf{UWS} --- underwater sensors on the seabed, measuring temperature (T),
      salinity (SAL), pressure (P) and current velocity (V).
\item \textbf{SUB} --- a submarine or AUV acting as a mobile hub that aggregates readings
      from several sensors.
\item \textbf{B} --- surface buoys relaying between the submarine and the satellite. Buoys
      carry GPS, so they contribute location as well as environmental data.
\item \textbf{SAT} --- satellites providing the long-distance link to shore.
\item \textbf{BS} --- the base station on land, which processes the data and acts as the
      trust anchor for revocation.
\end{itemize}

Crucially, the model builds in \textbf{redundancy}: each submarine keeps links to more than
one buoy, and each buoy can reach an alternate satellite. This is what makes the fallback
mechanism of Section~\ref{sec:fallback} possible.

\begin{figure}[H]
\centering
\begin{tikzpicture}[y=1.28cm,x=1.28cm]
% layer bands
\begin{scope}[on background layer]
\foreach \y/\lab in {0/{Layer 1 \textemdash\ UWS}, -1.5/{Layer 2 \textemdash\ SUB},
                     -3.0/{Layer 3 \textemdash\ BUOY}, -4.5/{Layer 4 \textemdash\ SAT},
                     -6.0/{Layer 5 \textemdash\ BS}}
  {\fill[mist] (-4.9,\y-0.45) rectangle (5.4,\y+0.45);
   \node[font=\tiny\sffamily,text=slate,anchor=west] at (-4.85,\y+0.62) {\lab};}
\end{scope}
% nodes
\node[node] (u1) at (-1.6, 0)    {U1};
\node[node] (u2) at ( 1.0, 0)    {U2};
\node[livenode] (s1) at (-0.3,-1.5) {S1};
\node[deadnode] (b1) at (-2.9,-3.0) {B1};
\node[livenode] (b2) at ( 2.0,-3.0) {B2};
\node[node,fill=mistdark] (sa1) at (-2.9,-4.5) {SAT1};
\node[node,fill=mistdark] (sa2) at ( 2.0,-4.5) {SAT2};
\node[node,fill=paleamber,draw=amber] (bs) at ( 0.0,-6.0) {BS};
% dead path
\draw[dead] (s1) -- (b1);
\draw[dead] (b1) -- (sa1);
\draw[flow,draw=slate,dashed] (sa1) -- (bs);
% active path
\draw[active] (u1) -- (s1);
\draw[active] (u2) -- (s1);
\draw[active] (s1) -- node[lblsm,fill=white,inner sep=1pt,pos=.55]{850\,m} (b2);
\draw[active] (b2) -- node[lblsm,fill=white,inner sep=1pt]{1000\,m} (sa2);
\draw[active] (sa2) -- node[lblsm,fill=white,inner sep=1pt,pos=.4]{500\,m} (bs);
\node[lblsm,anchor=east] at (-0.55,-0.78) {150\,m};
\node[lblsm,anchor=west] at (0.35,-0.78) {200\,m};
% annotations
\node[font=\scriptsize\bfseries,text=coral,anchor=east] at (-3.35,-3.0) {FAILED};
\node[box,anchor=west,align=left,fill=palekelp,draw=kelp,font=\tiny] at (2.9,-1.5)
  {\textbf{Active route}\\ U1/U2 $\to$ S1 $\to$ B2\\ $\to$ SAT2 $\to$ BS};
\node[box,anchor=west,align=left,fill=palecoral,draw=coral,font=\tiny] at (-4.85,-5.25)
  {\textbf{Dead branch}\\ B1 unresponsive,\\ SAT1 unreachable};
\end{tikzpicture}
\caption[UWC network model with fallback routing]{\textbf{The network model as we
instantiate it.} Eight nodes over five layers. Buoy B1 is forced to fail in the simulation
(dashed red), so the submarine reroutes through the pre-registered fallback buoy B2 and its
satellite SAT2 (solid green). Distances are the inter-node values used in our channel model
(Section~\ref{sec:thorp}).}
\label{fig:topology}
\end{figure}

\section{Overview of the three phases}

\begin{figure}[H]
\centering
\resizebox{\textwidth}{!}{%
\begin{tikzpicture}[x=1cm,y=1cm]
\node[box,minimum width=40mm,minimum height=17mm,fill=mist,draw=navy,align=left] (p1) at (-4.6,0)
  {\textbf{Phase 1 \textemdash\ Initialisation}\\[2pt]
   pick curve $\psi$, generator $\rho$\\
   $\PK_i \leftarrow \mathrm{rand}(1,n{-}1)$\\
   $\PU_i \leftarrow \PK_i \times \rho$\\
   $\ID_i \leftarrow H(\PU_i \cat E)$};
\node[box,minimum width=40mm,minimum height=17mm,fill=mist,draw=navy,align=left] (p2) at (0.4,0)
  {\textbf{Phase 2 \textemdash\ Registration}\\[2pt]
   $\RID_i \leftarrow H(\ID_i\cat\PU_i\cat\PK_i)$\\
   $\vartheta_i \leftarrow H(\RID_i\cat\PU_i)$\\
   peer recomputes $\vartheta_j$\\
   match $\Rightarrow$ add to set $R$};
\node[box,minimum width=40mm,minimum height=17mm,fill=mist,draw=navy,align=left] (p3) at (5.4,0)
  {\textbf{Phase 3 \textemdash\ Mutual auth.}\\[2pt]
   challenge with nonce\\
   response echoes nonce\\
   check $|\TS_2-\TS_1|<\Delta T$\\
   session $S_{i\leftrightarrow j}$ opens};
\draw[flow,line width=1pt] (p1) -- node[lblsm,above]{once, at} node[lblsm,below]{deployment} (p2);
\draw[flow,line width=1pt] (p2) -- node[lblsm,above]{once per} node[lblsm,below]{peer pair} (p3);
\node[box,fill=palekelp,draw=kelp,minimum width=30mm,font=\scriptsize] (d) at (5.4,-2.1)
  {then: encrypted data flows};
\draw[active] (p3) -- (d);
\draw[->,draw=slate,thick,dashed] (d.east) .. controls +(1.6,0) and +(1.6,-1.2) .. (p3.east)
  node[lblsm,pos=.5,right,align=left]{re-run per\\ hop, per\\ session};
\end{tikzpicture}}
\caption[The three protocol phases]{\textbf{The three phases.} Initialisation and
registration happen once; mutual authentication runs on every hop of every session and is
where the protocol's per-message cost lives.}
\label{fig:phases}
\end{figure}

\section{Phase 1: System initialisation}

Every entity $E \in \{\mathrm{UWS, SUB, B, SAT, BS}\}$ performs six steps:

\begin{enumerate}
\item Select the pentatope elliptic curve $\psi = X^3+Y^3+Z^3+W^3+T^3 = d\,XYZWT$ over
      $\mathbb{F}_p$ for a large prime $p$ and constant $d$.
\item Choose a generator point $\rho = (X_\rho,Y_\rho,Z_\rho,W_\rho,T_\rho)$ on $\psi$ of
      large prime order $n$.
\item Generate a private key $\PK_i \in \langle 1, n-1 \rangle$ uniformly at random.
\item Compute the public key $\PU_i = \PK_i \times \rho$, where $\times$ denotes scalar
      multiplication on $\psi$.
\item Derive the identifier by hashing the public key, $\ID_i = H(\PU_i \cat E)$. Because
      $H$ is one-way and collision-resistant, an identity cannot be forged without the
      corresponding key pair.
\item Publish $M_i = (\ID_i, \PU_i)$ to the entities it must talk to.
\end{enumerate}

\begin{keybox}
\textbf{Why derive the identity from the key.} If identities were assigned arbitrarily, an
attacker could claim any identity she liked. Binding $\ID_i = H(\PU_i \cat E)$ means that to
claim an identity you must possess a public key that hashes to it --- and to \emph{use} that
identity you must hold the matching private key. Impersonation therefore reduces to
inverting a hash or solving the ECDLP. This is the mechanism behind Proposition~1 in
Section~\ref{sec:informal}.
\end{keybox}

\section{Phase 2: Registration}

Registration binds an entity to the network and gives its peers a way to detect a tampered
or duplicate registration.

\begin{enumerate}
\item $E$ computes its registration identifier and a verification value:
      \begin{equation}
      \RID_i = H(\ID_i \cat \PU_i \cat \PK_i), \qquad
      \vartheta_i = H(\RID_i \cat \PU_i)
      \end{equation}
\item $E$ sends the registration message $M_i = \{\RID_i, \PU_i, \vartheta_i\}$ to the peer
      $j$.
\item $j$ extracts $\{\RID_i, \PU_i\}$ and independently recomputes
      $\vartheta_j = H(\RID_i \cat \PU_i)$. If $\vartheta_j \neq \vartheta_i$ the message
      has been altered in transit and is rejected as \textsc{invalid}.
\item If the message is valid but $\RID_i \in R$, the registration is a duplicate and is
      refused.
\item Otherwise $\RID_i$ is added to $R$ and a session $S_{i\leftrightarrow j}$ is created.
\end{enumerate}

Note that $\RID_i$ mixes in the \emph{private} key $\PK_i$, so no other party can compute
another node's $\RID$; but because $H$ is one-way, publishing $\RID_i$ does not leak
$\PK_i$.

\subsection{Storage cost of the registration set}

The paper is careful about memory on constrained nodes. The set $R$ is a hash table or Bloom
filter giving $O(1)$ lookup, and each entity stores only the subset relevant to its own
neighbourhood --- a sensor keeps the identifiers of its submarines, not of the whole
network. With 160-bit $\RID$ values and $k$ direct peers, storage is about $20k$ bytes; for
a sensor with fewer than ten peers this is under 250 bytes. Centralised nodes such as the BS
hold the full $O(N)$ record while constrained nodes hold an $O(k)$ partial view with
$k \ll N$.

\section{Phase 3: Mutual authentication}

This is the operational core. The chain from seabed to shore is completed in twelve
messages, in a pattern that repeats identically at each of the four hops:

\begin{itemize}
\item a \textbf{challenge} carrying the initiator's identity, a fresh nonce and a timestamp;
\item an \textbf{acknowledgement} carrying the responder's identity, the \emph{same} nonce
      echoed back, and a new timestamp;
\item a \textbf{data message} carrying the payload once the session is open.
\end{itemize}

Four hops $\times$ three messages $=$ twelve. Figure~\ref{fig:msc} is the complete sequence.

\begin{figure}[H]
\centering
\resizebox{\textwidth}{!}{%
\begin{tikzpicture}[x=1cm,y=1cm,font=\footnotesize]
% lifelines
\foreach \x/\nm in {0/UWS, 3.4/SUB, 6.8/B, 10.2/SAT, 13.6/BS}
 {\node[node,minimum width=16mm] at (\x,0) {\nm};
  \draw[draw=slate,dashed] (\x,-0.42) -- (\x,-13.9);}

%================= HOP 1 : UWS <-> SUB  (labels anchored WEST at x=0)
\begin{scope}[on background layer]\fill[mist] (-0.9,-1.15) rectangle (14.5,-3.85);\end{scope}
\node[font=\scriptsize\bfseries,text=teal,anchor=west] at (-0.85,-0.9) {HOP 1 \ \ UWS $\leftrightarrow$ SUB};
\node[lblsm,anchor=west] at (0,-1.35) {$M_1$ \ $\{\ID_{UWS},\, N_u,\, \TS_1\}_{pk(SUB)}$};
\draw[flow] (0,-1.62) -- (3.4,-1.62);
\node[lblsm,anchor=west] at (0,-2.05) {$M_2$ \ $\{\ID_{SUB},\, N_u,\, \TS_2\}_{pk(UWS)}$ \ \ \textit{\textcolor{teal}{same nonce echoed}}};
\draw[flow] (3.4,-2.32) -- (0,-2.32);
\node[box,fill=palekelp,draw=kelp,font=\tiny,anchor=west] at (0,-2.85)
  {$N_u$ matches \ \textbf{and} \ $|\TS_2-\TS_1|<\Delta T$ \ $\Rightarrow$ \ session $S_{UWS\leftrightarrow SUB}$ opens};
\node[lblsm,anchor=west] at (0,-3.35) {$M_3$ \ $\{\ID,\, \mathtt{Temp},\mathtt{Press},\mathtt{CurrVel},\mathtt{Sal},\, \TS_3\}_{pk(SUB)}$};
\draw[active] (0,-3.62) -- (3.4,-3.62);

%================= HOP 2 : SUB <-> B  (labels anchored WEST at x=3.4)
\node[font=\scriptsize\bfseries,text=teal,anchor=west] at (-0.85,-4.3) {HOP 2 \ \ SUB $\leftrightarrow$ B};
\node[lblsm,anchor=west] at (3.4,-4.75) {$M_4$ \ $\{\ID_{SUB},\, N_s,\, \TS_4\}_{pk(B)}$};
\draw[flow] (3.4,-5.02) -- (6.8,-5.02);
\node[lblsm,anchor=west] at (3.4,-5.45) {$M_5$ \ $\{\ID_{B},\, N_s,\, \TS_5\}_{pk(SUB)}$};
\draw[flow] (6.8,-5.72) -- (3.4,-5.72);
\node[box,fill=palekelp,draw=kelp,font=\tiny,anchor=west] at (3.4,-6.25)
  {$N_s$ matches \ \textbf{and} \ $|\TS_5-\TS_4|<\Delta T$ \ $\Rightarrow$ \ $S_{SUB\leftrightarrow B}$};
\node[lblsm,anchor=west] at (3.4,-6.75) {$M_6$ \ $\{\ID_{B},\, \mathtt{ENV},\, \mathtt{Alerts},\, \TS_6\}_{pk(B)}$};
\draw[active] (3.4,-7.02) -- (6.8,-7.02);

%================= HOP 3 : B <-> SAT  (labels anchored EAST at x=10.2)
\begin{scope}[on background layer]\fill[mist] (-0.9,-7.55) rectangle (14.5,-10.25);\end{scope}
\node[font=\scriptsize\bfseries,text=teal,anchor=west] at (-0.85,-7.75) {HOP 3 \ \ B $\leftrightarrow$ SAT};
\node[lblsm,anchor=east] at (10.2,-8.15) {$M_7$ \ $\{\ID_{B},\, N_b,\, \TS_7\}_{pk(SAT)}$};
\draw[flow] (6.8,-8.42) -- (10.2,-8.42);
\node[lblsm,anchor=east] at (10.2,-8.85) {$M_8$ \ $\{\ID_{SAT},\, N_b,\, \TS_8\}_{pk(B)}$};
\draw[flow] (10.2,-9.12) -- (6.8,-9.12);
\node[box,fill=palekelp,draw=kelp,font=\tiny,anchor=east] at (10.2,-9.62)
  {$N_b$ matches \ \textbf{and} \ $|\TS_8-\TS_7|<\Delta T$ \ $\Rightarrow$ \ $S_{B\leftrightarrow SAT}$};
\node[lblsm,anchor=east] at (10.2,-10.1) {$M_9$ \ $\{\ID,\, \mathtt{ENV},\, \mathtt{Alerts},\, \mathtt{GPS},\, \TS_9\}_{pk(SAT)}$};
\draw[active] (6.8,-10.37) -- (10.2,-10.37);

%================= HOP 4 : SAT <-> BS  (labels anchored EAST at x=13.6)
\node[font=\scriptsize\bfseries,text=teal,anchor=west] at (-0.85,-11.0) {HOP 4 \ \ SAT $\leftrightarrow$ BS};
\node[lblsm,anchor=east] at (13.6,-11.4) {$M_{10}$ \ $\{\ID_{SAT},\, N_{sat},\, \TS_{10}\}_{pk(BS)}$};
\draw[flow] (10.2,-11.67) -- (13.6,-11.67);
\node[lblsm,anchor=east] at (13.6,-12.1) {$M_{11}$ \ $\{\ID_{BS},\, N_{sat},\, \TS_{11}\}_{pk(SAT)}$};
\draw[flow] (13.6,-12.37) -- (10.2,-12.37);
\node[box,fill=palekelp,draw=kelp,font=\tiny,anchor=east] at (13.6,-12.87)
  {$N_{sat}$ matches \ \textbf{and} \ $|\TS_{11}-\TS_{10}|<\Delta T$ \ $\Rightarrow$ \ $S_{SAT\leftrightarrow BS}$};
\node[lblsm,anchor=east] at (13.6,-13.35) {$M_{12}$ \ $\{\ID,\, \mathtt{ENV},\, \mathtt{Alerts},\, \mathtt{GPS},\, \mathtt{LOGS},\, \TS_{12}\}_{pk(BS)}$};
\draw[active] (10.2,-13.62) -- (13.6,-13.62);

%================= legend
\draw[flow] (-0.85,-14.45) -- (0.35,-14.45);
\node[font=\tiny,anchor=west] at (0.45,-14.45) {authentication exchange};
\draw[active] (5.2,-14.45) -- (6.4,-14.45);
\node[font=\tiny,anchor=west] at (6.5,-14.45) {payload, sent only after the session opens};
\node[font=\tiny,anchor=west,text=slate] at (-0.85,-15.0)
  {Notation: $\{m\}_{pk(X)}$ means $m$ encrypted so that only $X$ can read it. \ $N_u,N_s,N_b,N_{sat}$ are the per-hop nonces \textemdash\ none of them crosses a hop boundary.};
\end{tikzpicture}}
\caption[Complete twelve-message authentication chain]{\textbf{The full protocol,
$M_1$ through $M_{12}$.} Each hop repeats the same three-message pattern: challenge,
acknowledgement echoing the nonce, then payload. Sessions are established
\emph{independently} per hop --- the sensor never talks to the base station directly, and no
nonce crosses a hop boundary. The payload grows as it climbs: sensor readings at hop~1, plus
alerts at hop~2, plus GPS at hop~3, plus logistics at hop~4.}
\label{fig:msc}
\end{figure}

\subsection{Steps in words}

\begin{enumerate}[label=\textbf{M\arabic*.},leftmargin=11mm]
\item UWS $\to$ SUB. The sensor encrypts $\langle \ID_{UWS} \cat \Nonce_{UWS} \cat \TS_1
      \rangle$ under the submarine's public key. Only the submarine, holding $\PK_{SUB}$,
      can decrypt --- so a successful decryption is already evidence the message was
      intended for it.
\item SUB $\to$ UWS. The submarine replies with its own identity, \emph{the same nonce}, and
      a fresh timestamp, encrypted under the sensor's public key. Echoing the nonce is the
      proof of liveness: only a party that could decrypt $M_1$ knows what to echo.
\item Session check. UWS verifies $\Nonce^{recv}_{SUB} = \Nonce^{sent}_{UWS}$ and
      $|\TS_2 - \TS_1| < \Delta T$. Both must hold. The session $S_{UWS \leftrightarrow SUB}$
      opens, and UWS sends the sensor payload
      $\langle$Temp, Press, CurrVel, Sal$\rangle$ as $M_3$. SUB aggregates it into ENV.
\item[\textbf{M4--M6.}] The identical pattern runs between SUB and B, with SUB's nonce. At
      $M_6$ the submarine forwards ENV plus \emph{Alerts} --- notifications of jamming,
      spoofing or eavesdropping detected lower down.
\item[\textbf{M7--M9.}] The pattern runs between B and SAT. At $M_9$ the buoy adds its
      \emph{GPS} coordinates to the payload.
\item[\textbf{M10--M12.}] The pattern runs between SAT and BS. At $M_{12}$ the satellite
      adds machine and logistics updates (\emph{LOGS}), and the base station receives the
      complete record.
\end{enumerate}

\begin{keybox}
\textbf{The central design decision.} Freshness is verified \emph{per hop}. Each hop
generates its own nonce and validates its own timestamp difference. Nothing propagates
end-to-end, and there is no global session state to track. This is what makes the protocol
viable underwater: a lost packet costs one hop's retry, not a whole-chain restart, and the
1.7\,s end-to-end propagation delay never enters a freshness check.
\end{keybox}

\section{A worked example: one message end to end}
\label{sec:worked}

Abstract notation hides how little actually happens per message. Figure~\ref{fig:pipeline}
traces $M_1$ from sensor U1 to submarine S1 using the real field sizes and real value shapes
our simulator produces.

\begin{figure}[H]
\centering
\resizebox{\textwidth}{!}{%
\begin{tikzpicture}[x=1cm,y=1cm,font=\scriptsize]
% ---- sender column
\node[font=\scriptsize\bfseries\sffamily,text=navy,anchor=west] at (-6.6,0.55) {SENDER \textemdash\ U1};
\node[box,anchor=west,align=left,minimum width=54mm,fill=mist] (s1) at (-6.6,0)
 {\textbf{1. Collect the three fields}\\
  $\ID_{U1}$ \ \ \texttt{3147de24777a6cae\ldots} \ (64 bits used)\\
  $\Nonce$ \ \ \ \texttt{48231} \ (64 bits, fresh)\\
  $\TS_1$ \ \ \ \ \ \texttt{1772451003.412} \ (8 bits)};
\node[box,anchor=west,align=left,minimum width=54mm,fill=mist] (s2) at (-6.6,-1.5)
 {\textbf{2. Concatenate with the separator}\\
  \texttt{"3147de24\ldots|48231|1772451003.412"}};
\node[box,anchor=west,align=left,minimum width=54mm,fill=mist,draw=teal] (s3) at (-6.6,-3.05)
 {\textbf{3. Derive the session key by ECDH}\\
  $P = a_{U1}\cdot \PU_{S1}$, \ \ $K=\mathrm{SHA256}(P.x)$\\
  $K=$ \texttt{9d2f7c0a\ldots} \ (256 bits)};
\node[box,anchor=west,align=left,minimum width=54mm,fill=mist,draw=teal] (s4) at (-6.6,-4.75)
 {\textbf{4. AES-256-GCM encrypt}\\
  IV $\leftarrow$ 12 random bytes\\
  $(C, \tau) \leftarrow \mathrm{AESGCM}_K(\mathrm{IV}, P)$};
\node[box,anchor=west,align=left,minimum width=54mm,fill=paleamber,draw=amber] (s5) at (-6.6,-6.35)
 {\textbf{5. Wire format transmitted}\\
  \texttt{[ IV 12 B ][ ciphertext 55 B ][ tag 16 B ]}\\
  total 83 bytes $=$ 664 bits};
\foreach \a/\b in {s1/s2, s2/s3, s3/s4, s4/s5} \draw[flow] (\a) -- (\b);

% ---- channel
\node[box,anchor=north,align=center,minimum width=30mm,minimum height=34mm,
      fill=white,draw=slate,dashed] (ch) at (0.35,0.3)
 {\textbf{\textcolor{slate}{ACOUSTIC}}\\ \textbf{\textcolor{slate}{CHANNEL}}\\[4pt]
  150\,m at 1500\,m/s\\ $=0.100$\,s propagation\\[3pt]
  $+$ Thorp absorption\\ $+$ multipath jitter\\[3pt]
  $\Rightarrow \approx 0.104$\,s\\[6pt]
  \textcolor{coral}{15\,\% chance the}\\ \textcolor{coral}{packet is lost}\\
  \textcolor{coral}{(up to 3 retries)}};
\draw[flow,line width=1.1pt] (-1.0,-3.2) -- (ch.west|-0,-3.2);
\draw[flow,line width=1.1pt] (ch.east|-0,-3.2) -- (2.7,-3.2);
\node[font=\tiny\ttfamily,text=coral,anchor=south] at (0.35,-4.6)
  {an eavesdropper here sees only random-looking bytes};

% ---- receiver column
\node[font=\scriptsize\bfseries\sffamily,text=navy,anchor=west] at (2.9,0.55) {RECEIVER \textemdash\ S1};
\node[box,anchor=west,align=left,minimum width=52mm,fill=mist,draw=teal] (r1) at (2.9,0)
 {\textbf{6. Derive the same key}\\
  $P' = b_{S1}\cdot \PU_{U1} = P$\ \ (ECDH commutes)\\
  $K' = K =$ \texttt{9d2f7c0a\ldots}};
\node[box,anchor=west,align=left,minimum width=52mm,fill=mist,draw=teal] (r2) at (2.9,-1.65)
 {\textbf{7. Split and verify the tag}\\
  IV $=$ first 12 B, \ tag $=$ last 16 B\\
  recompute $\tau'$; if $\tau' \neq \tau$ \textcolor{coral}{\textbf{abort}}};
\node[box,anchor=west,align=left,minimum width=52mm,fill=mist] (r3) at (2.9,-3.3)
 {\textbf{8. Decrypt and split on \texttt{|}}\\
  \texttt{p[0]}$=\ID$, \texttt{p[1]}$=\Nonce$, \texttt{p[2]}$=\TS$};
\node[box,anchor=west,align=left,minimum width=52mm,fill=mist,draw=amber] (r4) at (2.9,-4.95)
 {\textbf{9. Two freshness checks}\\
  \texttt{int(p[1]) == 48231} \ \ (nonce echo)\\
  \texttt{abs(now - float(p[2])) < 5.0}};
\node[box,anchor=west,align=left,minimum width=52mm,fill=palekelp,draw=kelp] (r5) at (2.9,-6.5)
 {\textbf{10. Both pass}\\
  charge 48.8\,\uJ, open $S_{U1\leftrightarrow S1}$,\\
  reply with $M_2$};
\foreach \a/\b in {r1/r2, r2/r3, r3/r4, r4/r5} \draw[flow] (\a) -- (\b);
\end{tikzpicture}}
\caption[One message traced through every stage]{\textbf{$M_1$ from U1 to S1, stage by
stage.} The left column is what the sender does, the right column what the receiver does,
and the centre is the physical channel. Steps~3 and~6 produce the same key without it ever
crossing the channel. Steps~4 and~7 are our AES-GCM substitution
(Section~\ref{sec:aes}); the base paper writes this stage as a public-key encryption under
$\PU_{SUB}$, which is cryptographically equivalent in effect and cheaper in practice. Note
that the tag is checked at step~7 \emph{before} any plaintext is trusted.}
\label{fig:pipeline}
\end{figure}

\subsection{Where the 2112 bits come from}

The paper's headline communication figure is 2112 bits for a complete authentication cycle.
Its accounting uses standardised field sizes: 64-bit identities, 64-bit nonces, 8-bit
timestamps, 160-bit hash outputs, and 512-bit representations for ECC points or AES keys.
Summing the fields transmitted across the login and authentication phases of all four hops
gives 2112 bits, against 3008--3216 bits for the five comparison schemes. At a nominal
modem rate of 10\,kbps, 2112 bits occupy
\begin{equation}
t_{tx} = \frac{2112 \text{ bits}}{10\,000 \text{ bits/s}} = 211\text{ ms},
\end{equation}
which is the reference line drawn on our delay plot in Chapter~\ref{ch:results}.

\section{Resilience through fallback paths}
\label{sec:fallback}

A relay in the sea is not a reliable component. Buoys are struck by shipping, drift out of
position, flood, or simply run down. The base paper's answer is \emph{dynamic reassignment
to fallback authentication peers}.

If a buoy $B_1$ stops responding, the submarine reinitiates mutual authentication with an
alternate buoy $B_2$ taken from its neighbour list. Likewise a buoy can route to a secondary
satellite when the primary fails or is delayed. The essential detail is that these fallback
entities are \textbf{pre-registered during system initialisation} and held in a trusted
cache. Rerouting therefore does not require a new registration exchange --- the public keys
and registration identifiers are already known --- and the same nonce-freshness and
timestamp bounds validate the rerouted session as would validate an original one.

\begin{figure}[H]
\centering
\resizebox{0.96\textwidth}{!}{%
\begin{tikzpicture}[x=1cm,y=1cm,font=\scriptsize]
%---------------- BEFORE
\node[font=\small\bfseries\sffamily,text=navy,anchor=west] at (0,1.5) {BEFORE \textemdash\ primary path healthy};
\node[node] (a1) at (0.8,0.6) {S1};
\node[node] (a2) at (3.2,0.6) {B1};
\node[node] (a3) at (5.6,0.6) {SAT1};
\node[node,fill=paleamber,draw=amber] (a4) at (8.0,0.6) {BS};
\node[node,fill=white,draw=slate] (a5) at (3.2,-0.9) {B2};
\node[node,fill=white,draw=slate] (a6) at (5.6,-0.9) {SAT2};
\draw[active] (a1) -- (a2); \draw[active] (a2) -- (a3); \draw[active] (a3) -- (a4);
\draw[draw=slate,dashed] (a1) -- (a5); \draw[draw=slate,dashed] (a5) -- (a6);
\draw[draw=slate,dashed] (a6) -- (a4);
\node[font=\tiny,text=slate,anchor=west] at (0.0,-1.7)
  {B2 and SAT2 are registered and cached, but idle};
%---------------- arrow
\node[font=\small\bfseries,text=coral] at (9.6,-0.15) {$\Longrightarrow$};
\node[font=\tiny,text=coral,align=center] at (9.6,-0.75) {B1 stops\\ responding};
%---------------- AFTER
\node[font=\small\bfseries\sffamily,text=kelp,anchor=west] at (11.0,1.5) {AFTER \textemdash\ rerouted, session preserved};
\node[node] (b1) at (11.8,0.6) {S1};
\node[deadnode] (b2) at (14.2,0.6) {B1};
\node[node,fill=white,draw=slate] (b3) at (16.6,0.6) {SAT1};
\node[node,fill=paleamber,draw=amber] (b4) at (19.0,0.6) {BS};
\node[livenode] (b5) at (14.2,-0.9) {B2};
\node[livenode] (b6) at (16.6,-0.9) {SAT2};
\draw[dead] (b1) -- (b2);
\draw[draw=slate,dashed] (b2) -- (b3); \draw[draw=slate,dashed] (b3) -- (b4);
\draw[active] (b1) -- (b5); \draw[active] (b5) -- (b6); \draw[active] (b6) -- (b4);
\node[font=\tiny,text=kelp,anchor=west,align=left] at (11.0,-1.7)
  {S1 selects the next active peer from its cache and runs $M_4$--$M_6$ with B2.\\
   The already-completed hop U1$\leftrightarrow$S1 is \emph{not} torn down.};
\end{tikzpicture}}
\caption[Dynamic fallback rerouting]{\textbf{Fallback in action.} Because $B_2$ was
registered at initialisation, the submarine can authenticate with it immediately; only the
three messages of the failed hop are repeated. Nothing upstream is affected, which is the
direct benefit of authenticating per hop rather than end to end.}
\label{fig:fallback}
\end{figure}

The paper reports that a new session can be established in under 1\,ms thanks to the
lightweight cryptographic operations and pre-distributed public keys, so a node that drifts
into a new submarine's range can simply re-run $M_1$--$M_3$ rather than waiting on a
recovery procedure. Ongoing sessions remain valid as long as latency stays within
$\Delta T$ and loss stays below threshold.

\section{Long-term key management}

Three provisions in the paper address cryptographic sustainability over a deployment
lifetime.

\begin{enumerate}
\item \textbf{Secure private-key storage.} Private keys are held in tamper-resistant
      hardware --- a secure element such as the Microchip ATECC608A, or a TPM --- where cost
      and power allow. Where they do not, the paper falls back on firmware-level obfuscation
      and memory isolation to raise the cost of physical extraction.
\item \textbf{Key revocation.} The base station periodically issues a signed revocation
      list naming compromised or expired nodes. It is broadcast to the network, and every
      entity checks incoming messages against it before initiating or answering an
      authentication. The list itself is authenticated with a management key held by all
      trusted nodes.
\item \textbf{Initial public-key distribution.} At deployment each node uploads its public
      key to the base station over a trusted link --- a physical interface, or a short-range
      line-of-sight acoustic channel. The key is bound to the identity by
      \begin{equation}
      \ID_i = H(\PU_i \cat E \cat \mathrm{DeviceSerial})
      \end{equation}
      which prevents identity spoofing, and the keys are then disseminated to the other
      entities so mutual authentication can proceed.
\end{enumerate}
